% SOURCE_AUTHORITY: sga5_fr_workpass.tex, lines 7332--7353 (LF-normalized unit).
% SOURCE_SHA256: 791F4EFFC5E02832D5D77ED03518C8156D6F07E4C8238B03545DB93D883FBB28
\subsubsection*{4.1.4}
Se $X$ é um objeto de \(P\) que verifica (ML), deduz-se da expressão do graduado recordada em (4.1.1) que, para todo inteiro \(p\), o sistema projetivo
\[
\gr^p(X)=\gr^p(X_n)_{n\in\mathbb N}
\]
É essencialmente constante.

Seu limite projetivo, que existe sem hipótese sobre a categoria abeliana \(\mathcal C\), será chamado \(p\)-ésima componente do graduado estrito de \(X\), e será denotado por
\[
\grs^p(X).
\]

\subsection*{4.2. Sistemas projetivos adaptados a \(J\) e lema de Artin-Rees}

São dados um anel comutativo \(A\), um ideal \(J\) de \(A\) e uma \(A\)-categoria abeliana \(\mathcal C\).

\begin{statement}{Definição 4.2.1}
Um objeto \(X\) da categoria \(P=\underline{\Hom}(\mathbb N^\circ,\mathcal C)\) diz-se adaptado a \(J\) se sua filtração \((F^pX)_{p\in\mathbb N}\) de sistema projetivo (4.1.1) é compatível com a filtração \(J\)-ádica de \(X\), isto é, se para todo par \((p,q)\) de inteiros positivos ou nulos se verifica a inclusão que aparece a seguir:
\[
J^pF^q(X)\subset F^{p+q}(X).
\]
\end{statement}
